\scanpage{50}

\deleted{\unsure{On considère} \illegible{1}}

\uline{Avec} cor. \unsure{3} du th. de comparaison affine

Supposons que $\unsure{\mathfrak F'|\mathcal F'}$ satisfasse a) b) c')
[c') est que $\unsure{\mathfrak F'|\mathcal F'}$ s'étende en un
$\unsure{\mathfrak F|\mathcal F}$ coh. sur $\hat X$], \unsure{je}
dis que $\unsure{\mathfrak F'|\mathcal F'}$ provient d'un $F'$ satisfaisant la
condition $\operatorname{prof}F'_x\ge2$ si $x\in T_1$.
En effet,
il provient déjà d'un $F'$ coh. \deleted{\unsure{tel que}}\inserted{où}
$t$ est $F'$-régulier, il faut au besoin modifier $F'$.

\deleted{Soit
$Z'$ l'ens. des pts de $\unsure{X'}$ tels que
$\operatorname{prof}_{x'}(F')\leq1$,
c'est une partie \deleted{constructible \unsure{et} \illegible{3}}
\inserted{fermée de $X'-Y'$}.
On peut déjà \unsure{supposer} en
\unsure{effet} que pour $\forall$ $x'\in\operatorname{Ass}F'$, on ait
\deleted{$\unsure{x'\notin Y'}$ $\unsure{\dim\overline{x'}\in T_1}$}
\unsure{i.e.} $\overline{x'}\cap Y'\ne\varnothing$, \unsure{quitte à}
\unsure{limiter} $\uline{F'}$ \unsure{à un ouvert convenable}.
\unsure{Considérons} l'ensemble $U$ des
$x'\in X'$ tels que $\operatorname{prof}_{x'}(F')\ge2$.}

\margin{adhérence $\overline{x'}$ dans $X'$}

\margin{On suppose $X$ \unsure{localisé}
autour de $T$, et
$X-T=X'$ \deleted{\unsure{partie}}\inserted{un}
\deleted{\illegible{2}}
\uline{Jacobson} ``\unsure{sympathique}''}

Si $x'$ est
un pt \uline{fermé} de \deleted{$\unsure{X}$} $Y'$, on a : par hyp.
$x'\in U$, c'est donc vrai pour $\forall$ pt de $Y'$,
donc $U$ est un \uline{voisinage} de $Y'$ dans $X'$.

\deleted{Je dis que $x'\in X'-U$ implique que}
Les pts de $X'-U$ sont donc \uline{fermés} dans $X'$
[et sont en fait dans $T_1$] -- ils sont à
fortiori en nbre fini. Soit $x'\in X'-U$, et
\uline{soit} \deleted{$x'_1\in\operatorname{Ass}F'$ \unsure{tel}}
$x'_1\in U$ tel que $x'_1\in\operatorname{Ass}F'$,
$\overline{x'_1}\ni x'$, je dis que \unsure{il n'est pas possible que}
$x'_1\uline{\ne}x'$ : sinon en effet, $\overline{x'_1}$ serait de dim $1$,

\margin{49}
